% ============================================================
% Come si usano, argomento 24: minimi quadrati e stima ricorsiva
% ============================================================
\argomento{Minimi quadrati e stima ricorsiva}{Formule a p.~\pageref*{f:modellolineare}}

\begin{scelta}{Come scegliere la formula}
\renewcommand{\arraystretch}{1.25}
\begin{tabular}{@{}P{0.55\linewidth}P{0.39\linewidth}@{}}
  \toprule
  \textbf{Se il testo\ldots} & \textbf{si usa\ldots} \\
  \midrule
  dà dati misurati, senza probabilità, e chiede il modello lineare o la retta che li descrive meglio &
    \vv{modellolineare} \vv{minimiquadrati} \\
  aggiunge un dato e chiede di aggiornare la stima senza rifare i conti & \vv{rls} \\
  chiede l'inversa di una matrice a cui si somma un termine $\vett{u}\vett{v}^T$ & \vv{shermanmorrison} \\
  \bottomrule
\end{tabular}

Per una retta $\theta = a_1 + a_2 t$ bastano cinque somme: $p$, $\sum t_i$, $\sum t_i^2$, $\sum \theta_i$, $\sum t_i
\theta_i$. Poi si inverte una matrice $2 \times 2$ (\vv{inversadet}).
\end{scelta}

\scheda{modellolineare}{
  \usoquando{Si vuole spiegare una grandezza misurata con una combinazione lineare di alcune caratteristiche: il prezzo
    di una casa da superficie, stanze e piano.}
  \usocome{Ogni dato diventa una colonna di $\matr{X}$ e il valore misurato un elemento di $\vett{y}$; per avere anche
    una costante si aggiunge a ogni $\vett{x}(i)$ una componente uguale a 1.}
  \usoesempio{Retta $\theta = a_1 + a_2 t$: $\vett{x}(i) = [1, t_i]^T$, quindi $n = 2$; con sei punti $\matr{X}$ è
    $2 \times 6$.}
}

\scheda{minimiquadrati}{
  \usoquando{Si cerca il modello lineare che rende minima la somma dei quadrati degli errori sui dati.}
  \usocome{Si calcolano $\matr{X}\matr{X}^T$ e $\matr{X}\vett{y}$ come somme sui dati e si risolvono le equazioni
    normali.}
  \usoattenzione{Se una caratteristica è combinazione lineare delle altre, $\matr{X}\matr{X}^T$ non è invertibile.}
  \usoesempio{Punti $t = 1, \ldots, 6$ e $\theta = 2{,}1;\ 3{,}9;\ 6{,}2;\ 7{,}8;\ 10{,}1;\ 12{,}0$:
    $\matr{X}\matr{X}^T = \begin{bmatrix} 6 & 21 \\ 21 & 91 \end{bmatrix}$, $\matr{X}\vett{y} = [42{,}1;\
    182{,}2]^T$, determinante 105, retta $\theta \simeq 0{,}047 + 1{,}991\, t$, somma dei quadrati degli errori
    $0{,}107$.}
}

\scheda{shermanmorrison}{
  \usoquando{Si conosce già l'inversa di $\matr{R}$ e alla matrice si aggiunge un termine $\vett{u}\vett{v}^T$, come
    quando arriva un dato nuovo nei minimi quadrati.}
  \usocome{Bastano il prodotto $\matr{R}^{-1}\vett{u}$ e un numero al denominatore: nessuna nuova inversione.}
  \usoesempio{$\matr{R} = \matr{I}$, $\vett{u} = \vett{v} = [1, 1]^T$: $\left( \matr{I} + \vett{u}\vett{u}^T
    \right)^{-1} = \matr{I} - \frac{1}{3}\begin{bmatrix} 1 & 1 \\ 1 & 1 \end{bmatrix} = \frac{1}{3} \begin{bmatrix} 2
    & -1 \\ -1 & 2 \end{bmatrix}$, l'inversa di $\begin{bmatrix} 2 & 1 \\ 1 & 2 \end{bmatrix}$.}
}

\scheda{rls}{
  \usoquando{I dati arrivano uno alla volta e si vuole aggiornare la stima a costo fisso, senza invertire ogni volta
    una matrice.}
  \usocome{Si calcolano $\matr{R}^{-1}(p)\vett{x}$ e $K$, si corregge la stima in proporzione all'errore che il vecchio
    modello fa sul dato nuovo, e si aggiorna $\matr{R}^{-1}$.}
  \usoesempio{Ai sei punti si aggiunge $(7;\ 14{,}1)$: la retta prevedeva $13{,}99$, l'errore è $0{,}11$, $K =
    \frac{91}{105} \simeq 0{,}87$. La stima diventa $[0{,}014;\ 2{,}004]^T$, la stessa che si ottiene da capo con sette
    punti.}
}
